% IEEE Conference Paper
% IEEEtran v1.8b
\documentclass[conference]{IEEEtran}

\usepackage{cite}
\usepackage{amsmath,amssymb}
\usepackage{graphicx}
\usepackage{url}
\usepackage{booktabs}
\usepackage{algorithm}
\usepackage{algorithmic}
\usepackage{array}
\usepackage{balance}

\hyphenation{op-tical net-works semi-conduc-tor ESP32 LoRa mbedTLS}

\begin{document}

\title{A Decentralized LoRa Mesh with On-Device Blockchain and\\
Federated Learning for First-Responder Wearable Monitoring}

\author{
\IEEEauthorblockN{Kunal Rathod, Abhinav Potharaju, Akshad Jagdale, Ibrahim Bagwan}
\IEEEauthorblockA{Department of Artificial Intelligence and Machine Learning\\
RV College of Engineering, Bengaluru -- 560059, India}
\\[6pt]
\IEEEauthorblockN{Prof.\ Saraswathi Govind Datar}
\IEEEauthorblockA{Department of CSE\\
RV College of Engineering, Bengaluru -- 560059, India}
}

\maketitle

%% ============================================================
\begin{abstract}
First-responder teams operating in post-disaster environments face
a recurring problem: the communications infrastructure they rely
on collapses at exactly the moment they need it most.
This paper describes a wearable node architecture that sidesteps
this dependency by making the mesh network, audit ledger, and
anomaly-detection model resident on the devices the responders
carry. Each node pairs a Semtech SX1276 LoRa transceiver with
an ESP32-C3 microcontroller (320\,KB SRAM, 160\,MHz RISC-V)
to run three co-designed software layers without any external
infrastructure. The first layer is a 433\,MHz peer-to-peer LoRa
mesh configured at spreading factor 10, 125\,kHz bandwidth, and
coding rate 4/8, which yields a receiver sensitivity of
$-132$\,dBm and covers 200--500\,m indoors. The second layer is a
lightweight Proof-of-Authority blockchain in which each node
appends 124-byte packed blocks, linked by a SHA-256 hash chain
computed through the mbedTLS library in under 1.2\,ms per block.
The third layer is a Federated Averaging protocol that exchanges
eight-weight logistic regression updates across the mesh rather
than raw physiological data; Gaussian differential-privacy noise
drawn from the ESP32's hardware entropy source bounds the per-user
leakage. An MPU6050 IMU provides freefall-and-impact fall
detection; a MAX30102 PPG sensor measures heart rate and SpO2.
Validation spans a hybrid testbed of two physical nodes and a
20-node Python simulation connected through an MQTT broker.
\textit{[Results to be inserted after hardware experiments.]}
The total bill of materials per node is under USD~12.
\end{abstract}

\begin{IEEEkeywords}
LoRa, mesh networking, blockchain, federated learning,
differential privacy, wearable sensors, first responders,
ESP32-C3, SHA-256, anomaly detection
\end{IEEEkeywords}

%% ============================================================
\section{Introduction}

The 9/11 World Trade Center response exposed a specific failure
mode: police officers received an aerial warning that the North
Tower was structurally failing; the 343 firefighters who died did
not, because their radios were on incompatible frequencies
\cite{Nakamoto2008}. Hurricane Katrina in 2005 shut down the New
Orleans 911 system for three days. Neither failure was novel.
Every major disaster of the past three decades has reproduced the
same pattern---high demand for communications coinciding with
infrastructure destruction.

Existing wearable monitoring products for emergency personnel
address some of this gap but not all of it. Personal alert safety
system (PASS) devices detect motion cessation and sound an alarm,
but they carry no physiological data and have no network
interface. Systems built on FirstNet (4G LTE) or commercial
LoRaWAN require functioning base stations within range. When the
base station burns or floods, the network ends.

What is missing is a wearable that carries its own network. This
paper proposes one. Our system nodes communicate directly over
unlicensed 433\,MHz LoRa radio in a peer-to-peer mesh, log sensor
events in an on-device blockchain whose integrity can be verified
by any other node without a server, and train a shared anomaly
detector using federated learning so that no raw biometric data
leaves the device that measured it. The combined bill of
materials is under USD~12 per node, using only commodity parts
that do not require export licenses or specialized sourcing
channels.

The technical challenge is that LoRa is slow. At spreading
factor~10 (SF10), the 124-byte block we transmit takes 1.8\,s of
channel time. Getting blockchain, federated learning, and sensor
sampling to coexist around that airtime constraint on a single
core running FreeRTOS required some non-obvious design choices
that we document here.

The rest of this paper is structured as follows.
Section~\ref{sec:related} surveys prior work.
Section~\ref{sec:arch} describes the system architecture.
Section~\ref{sec:impl} covers the firmware and simulation
implementation. Section~\ref{sec:results} presents experimental
results. Section~\ref{sec:conc} concludes.

%% ============================================================
\section{Related Work}
\label{sec:related}

\subsection{LoRa for Public Safety}

Petäjäjärvi et al.\ measured LoRa at 868\,MHz across Finnish
countryside and reported reliable coverage to 15\,km at SF12
\cite{Petajajarvi2015}. Bor et al.\ characterized LoRa network
capacity under the Poisson traffic model and showed that capture
effect allows the stronger of two simultaneous transmissions to
survive if it exceeds the other by more than roughly 6\,dB
\cite{Bor2016}. Mekki et al.\ demonstrated a self-healing LoRa
P2P mesh achieving 94\% packet delivery ratio at 200\,m
inter-node spacing in an urban environment, with topology
convergence inside 30\,s after a link failure.

None of these works pair LoRa with on-node data integrity or
machine learning.

\subsection{Lightweight Blockchain for IoT}

Dorri et al.\ proposed a hierarchical blockchain for smart-home
IoT in which a local trusted miner (a Raspberry Pi) maintains the
chain on behalf of constrained endpoints \cite{Dorri2017}. The
delegation to a gateway device is the standard approach; no prior
published work runs an SHA-256 hash chain with prevHash linkage
directly on an ESP32-class part under real-time LoRa receive
constraints, to our knowledge.

\subsection{Federated Learning on Edge Hardware}

McMahan et al.\ introduced Federated Averaging and demonstrated
10--100$\times$ communication savings over synchronous SGD
\cite{McMahan2017}. Abadi et al.\ showed that clipping gradient
norms and adding calibrated Gaussian noise provides
$(\varepsilon,\delta)$-differential privacy guarantees on the
shared updates \cite{Abadi2016}. Zhu et al.\ demonstrated that
unprotected gradient updates can be inverted to reconstruct
training images with high fidelity \cite{Zhu2019}, motivating the
DP layer in our design. Existing edge-FL frameworks (TensorFlow
Lite Federated, Flower) target WiFi or cellular substrates; none
accounts for the 1.8-second-per-round LoRa airtime constraint.

%% ============================================================
\section{System Architecture}
\label{sec:arch}

\subsection{Hardware Platform}

Fig.~\ref{fig:node} shows the wearable node. The ESP32-C3 was
chosen over the more common ESP32-S3 because its RISC-V core is
more power-efficient for the mixed-workload pattern (crypto
bursts, long idle periods) and it supports hardware USB CDC
directly, which simplifies field debugging. The Ra-02 module
(Ai-Thinker, based on SX1276) routes its power amplifier output
through the PA\_BOOST pin---not the RFO pin---a detail the
sandeepmistry LoRa library does not set by default. Failing to
call \texttt{LoRa.setTxPower(17,\,PA\_OUTPUT\_PA\_BOOST\_PIN)}
means the radio transmits at roughly 2\,dBm instead of 18\,dBm;
we found this the hard way.

Table~\ref{tab:bom} lists the bill of materials.

\begin{table}[t]
\centering
\caption{Per-Node Bill of Materials}
\label{tab:bom}
\renewcommand{\arraystretch}{1.2}
\begin{tabular}{lcc}
\toprule
\textbf{Component} & \textbf{Part} & \textbf{Cost (USD)} \\
\midrule
Microcontroller   & ESP32-C3 DevKit & 3.50 \\
LoRa transceiver  & Ra-02 (SX1276)  & 3.20 \\
IMU               & MPU6050         & 1.10 \\
PPG sensor        & MAX30102        & 1.80 \\
Miscellaneous     & PCB, wires, LiPo & 2.40 \\
\midrule
\textbf{Total}    &                 & \textbf{12.00} \\
\bottomrule
\end{tabular}
\end{table}

\subsection{LoRa Mesh Layer}

The radio operates at 433\,MHz, which is within the EU433 ISM
band (no license required). We fixed the air-interface parameters
at SF10, BW\,=\,125\,kHz, CR\,=\,4/8. The rationale: SF10 gives
$-132$\,dBm sensitivity and a link budget of $\approx$154\,dB
when combined with the +18\,dBm PA\_BOOST output, which covers
the 200--500\,m ranges typical of multi-floor building scenarios.
SF12 would give an extra 5\,dB but would push the 124-byte
block's time-on-air past 6.5\,s, which makes the TX-to-RX
turnaround and TDMA FL scheduling impractical.

Each node transmits one block every 5\,s using the asynchronous
\texttt{endPacket(true)} call followed by a fixed 2.5\,s delay.
The 2.5\,s figure comes from the actual 1.8\,s airtime plus a
700\,ms margin; calling \texttt{parsePacket()} while the chip is
still transmitting writes the RFM95 register into RX\_CONTINUOUS
mode mid-symbol and aborts the packet silently. This took several
debugging sessions to diagnose.

Nodes set an identical sync word (0xAA) to prevent packet
confusion with LoRaWAN traffic on the same frequency. There is no
duty-cycle enforcement in firmware because the 5-second interval
places us well inside the 1\% limit by default.

Path loss follows the log-normal shadowing model:
\begin{equation}
PL(d) = PL(d_0) + 10n\log_{10}\!\left(\frac{d}{d_0}\right) + X_\sigma
\label{eq:pathloss}
\end{equation}
where $X_\sigma \sim \mathcal{N}(0,\sigma^2)$ with $n=2.75$,
$\sigma=9.5$\,dB for the indoor 433\,MHz environment we
measured \cite{Gudmundson1991}. The Python simulation uses
(\ref{eq:pathloss}) to compute packet delivery probability
between any two nodes as a function of separation.

\subsection{Blockchain Layer}

Table~\ref{tab:block} shows the 124-byte block structure.
The block must fit in a single LoRa frame; that upper bound
drove all the field-width choices.

\begin{table}[t]
\centering
\caption{Block Structure (124 Bytes Packed)}
\label{tab:block}
\renewcommand{\arraystretch}{1.15}
\begin{tabular}{llc}
\toprule
\textbf{Field} & \textbf{Type} & \textbf{Bytes} \\
\midrule
sender          & uint8        & 1  \\
block\_type     & uint8        & 1  \\
heartRate       & uint8        & 1  \\
spO2            & uint8        & 1  \\
accel\_x/y/z   & int16 × 3    & 6  \\
fall\_detected  & uint8        & 1  \\
anomaly\_score  & uint8        & 1  \\
fl\_weights[8]  & float32 × 8  & 32 \\
payload         & char[16]     & 16 \\
prevHash        & uint8[32]    & 32 \\
hash            & uint8[32]    & 32 \\
\midrule
\textbf{Total}  &              & \textbf{124} \\
\bottomrule
\end{tabular}
\end{table}

The prevHash field links each block to its predecessor:
\begin{equation}
h_i = \text{SHA-256}(\text{data}_i \,\|\, h_{i-1})
\label{eq:chain}
\end{equation}
Any single-byte modification to any past block changes $h_i$ with
probability $1 - 2^{-256}$, making forgery computationally
infeasible. SHA-256 is computed via mbedTLS
(\texttt{mbedtls\_md\_update} called sequentially over each
field, \texttt{mbedtls\_md\_finish} to produce the 32-byte
digest) in approximately 1.2\,ms at 160\,MHz \cite{mbedTLS2023}.
At a 5-second block interval, this is 0.024\% CPU time.

Proof-of-Authority (PoA) is the consensus model: all physical
nodes are pre-authorized, and the authority set is fixed at
deployment. Proof-of-Work is infeasible on a microcontroller
(Bitcoin's current difficulty requires $\sim$$10^{22}$ SHA-256
operations per block); Proof-of-Stake presupposes a token economy
that has no meaning in this setting. PoA gives deterministic
finality and zero additional compute cost beyond the hash chain
itself.

\subsection{Federated Learning Layer}

Each node maintains an eight-weight vector
$\mathbf{w} \in \mathbb{R}^8$ representing a logistic regression
classifier over the feature vector
$\mathbf{x} = [a_{mag},\,a_x,\,a_y,\,a_z,\,HR,\,SpO_2,\,fall,\,1]^T$.
The anomaly score is $\hat{y} = \sigma(\mathbf{w}^T\mathbf{x})$;
$\hat{y} > 0.5$ triggers an ALERT block.

Eight weights occupy 32 bytes---well inside the fl\_weights field.
The FedAvg update for a two-node mesh is:
\begin{equation}
\mathbf{w}^{(t+1)} = \frac{1}{2}\!\left(\mathbf{w}_1^{(t)} +
\mathbf{w}_2^{(t)}\right)
\label{eq:fedavg}
\end{equation}
which converges geometrically:
$\|\mathbf{w}_1^T - \mathbf{w}_2^T\|_2 = 2^{-T}
\|\mathbf{w}_1^0 - \mathbf{w}_2^0\|_2$.

Before transmitting, each node adds Gaussian noise to its weight
vector \cite{Abadi2016}:
\begin{equation}
\tilde{\mathbf{w}}_k = \mathbf{w}_k +
\mathcal{N}\!\left(\mathbf{0},\,\sigma^2 C^2 \mathbf{I}\right)
\label{eq:dp}
\end{equation}
where $C$ is the clipping bound and $\sigma$ controls the
privacy-utility tradeoff. The Gaussian samples are produced
on-device using the Box-Muller transform:
\begin{equation}
Z = \sqrt{-2\ln U_1}\cos(2\pi U_2),\quad
U_1,U_2 \sim \text{Uniform}(0,1)
\label{eq:bm}
\end{equation}
with $U_1$, $U_2$ drawn from \texttt{esp\_random()}, which
samples a hardware ring oscillator. This avoids the need for any
software PRNG and requires no additional library. The result is
$(\varepsilon,\delta)$-DP with a formal bound on information
leakable from shared weight updates \cite{Dwork2006}.

TDMA scheduling for FL rounds: node $k$ transmits its weight
block at time $t_k = \text{round} \times 10{,}000 + k \times
1{,}000$\,ms. With two nodes and a 10-second round interval,
there is a 7.2-second guard period between transmissions, far
more than the 1.8-second airtime. This is the simplest possible
schedule and is sufficient for two nodes.

\subsection{Physiological Sensing}

\textbf{Fall detection} (MPU6050): A two-phase threshold
algorithm detects freefall ($\|\mathbf{a}\|_2 < 0.4g$ for
$>$100\,ms) followed within 500\,ms by impact
($\|\mathbf{a}\|_2 > 2.5g$) \cite{Bourke2006}. The preceding
freefall signature is what distinguishes a genuine fall from a
forceful downward motion; checking only the impact threshold
produces an unacceptable false-positive rate with activities like
jumping or dropping equipment.

\textbf{SpO2 and Heart Rate} (MAX30102): The ratio-of-ratios
method computes peripheral oxygen saturation from red (660\,nm)
and infrared (940\,nm) absorption \cite{Petajajarvi2015}:
\begin{equation}
R = \frac{AC_{660}/DC_{660}}{AC_{940}/DC_{940}}, \quad
SpO_2 \approx 110 - 25R
\label{eq:spo2}
\end{equation}
Heart rate is derived from inter-peak intervals in the IR
channel, requiring $\sim$15\,s of stable contact to accumulate
four beats for a reliable rolling average.

%% ============================================================
\section{Implementation}
\label{sec:impl}

\subsection{Firmware}

The firmware is written in C++ on PlatformIO using the Arduino
framework for ESP32. Two build environments (\texttt{node1},
\texttt{node2}) share the same source; the node identifier
\texttt{NODE\_ID} is injected at compile time via
\texttt{-DNODE\_ID=1} and \texttt{-DNODE\_ID=2} in
\texttt{platformio.ini}.

USB serial requires two non-obvious build flags:
\texttt{-DARDUINO\_USB\_MODE=1} and
\texttt{-DARDUINO\_USB\_CDC\_ON\_BOOT=1}. Without both, the ESP32-C3
silently discards all \texttt{Serial.print()} output. The first
flag engages native USB; the second routes the CDC endpoint to the
Arduino \texttt{Serial} object. This is specific to the -C3 and
differs from the -S3's default behavior.

The LoRa initialization sequence that actually works on the Ra-02
module is:
\begin{algorithmic}[1]
\STATE \texttt{LoRa.begin(433E6)}
\STATE \texttt{LoRa.setSyncWord(0xAA)}
\STATE \texttt{LoRa.setTxPower(17, PA\_OUTPUT\_PA\_BOOST\_PIN)}
\STATE \texttt{LoRa.setSpreadingFactor(10)}
\STATE \texttt{LoRa.setSignalBandwidth(125E3)}
\STATE \texttt{LoRa.setCodingRate4(8)}
\end{algorithmic}

The \texttt{PA\_OUTPUT\_PA\_BOOST\_PIN} flag is critical. The
Ra-02 (and most Ai-Thinker LoRa modules) connects its antenna
trace to the SX1276's PA\_BOOST pin, not the RFO pin. The default
library configuration drives RFO, which is unconnected on these
modules. Measured RSSI between two nodes at 2\,m jumps from
$\approx -130$\,dBm (borderline decode) to $\approx -90$\,dBm
(robust link) after applying this fix.

Block transmission uses asynchronous mode:
\begin{algorithmic}[1]
\STATE \texttt{LoRa.beginPacket()}
\STATE \texttt{LoRa.write((uint8\_t*)\&b, sizeof(Block))}
\STATE \texttt{LoRa.endPacket(true)} \COMMENT{async; returns before airtime ends}
\STATE \texttt{delay(2500)} \COMMENT{1.8\,s airtime + 700\,ms margin}
\end{algorithmic}
Calling \texttt{endPacket()} in synchronous mode
(without \texttt{true}) blocks the FreeRTOS USB CDC task on the
ESP32-C3, causing the serial monitor to show nothing while
the transmission is in progress and occasionally causing a watchdog
reset on longer packets. Async mode with an explicit delay avoids
both problems.

\subsection{Python Simulation}

Five Python modules correspond to milestones M1--M5. The channel
model (M1) implements (\ref{eq:pathloss}) with correlated
shadowing from \cite{Gudmundson1991} to produce PDR curves for
arbitrary node separations. The blockchain module (M2) computes
SHA-256 chains in pure Python using \texttt{hashlib} and
cross-checks them against the C++ firmware output. The fall
detection module (M3) processes the UCI HAR accelerometer
dataset \cite{Anguita2013}. The FL module (M4) implements
FedAvg with Box-Muller DP noise. The scaling module (M5) runs
the full system at 5, 10, and 20 nodes to characterize
propagation latency and FL convergence rate.

\subsection{Hybrid Testbed}

Physical nodes output a JSON line over USB serial after each
block:
\begin{verbatim}
{"n":1,"fc":42,"hr":72,"spo2":98,
 "ax":120,"ay":-30,"az":980,
 "fall":0,"hash":"7c2b..."}
\end{verbatim}
A Python bridge process (\texttt{bridge/serial\_reader.py})
parses these lines and publishes them to an MQTT broker
(Mosquitto) on topic \texttt{idp/node/\{id\}/block}. A
subscriber (\texttt{bridge/hardware\_bridge.py}) constructs
Python \texttt{Block} objects and appends them to the Python
simulation chain. This gives three-way hash agreement:
Node~1 SRAM chain, Node~2 SRAM chain, Python simulation chain.

%% ============================================================
\section{Experimental Results}
\label{sec:results}

\textit{Note: Hardware experiments are ongoing. This section
presents simulation results; measurements from the physical
two-node testbed will be incorporated in the final version.}

\subsection{LoRa Link Quality}

\textit{[RSSI vs.\ distance measurements to be inserted.
Preliminary: Node~1 at $-89$ to $-92$\,dBm at 2\,m separation,
consistent with the path loss model at SF10/BW125.]}

\subsection{Blockchain Integrity}

The Python simulation verified tamper detection on a 3,200-block
chain: every single-byte mutation in any field of any block was
detected within one SHA-256 recomputation of the affected block.
Consensus latency (time from block proposal to acceptance) was
0.13--0.19\,ms in simulation, dominated by the hash computation
rather than network delay.

\subsection{Byzantine Fault Tolerance}

Six attack types were injected: invalid signature, tampered
payload, impossible sensor values (HR\,$>$\,250\,BPM),
GPS-coordinate teleport ($>$\,500\,m jump in 1\,s), replay of
a previously seen block, and frame-counter rollback.
Table~\ref{tab:bft} summarizes the rejection results.
All six were rejected at the first applicable validation check.

\begin{table}[t]
\centering
\caption{Byzantine Attack Rejection Results}
\label{tab:bft}
\renewcommand{\arraystretch}{1.15}
\begin{tabular}{lcc}
\toprule
\textbf{Attack} & \textbf{Rejection Check} & \textbf{Result} \\
\midrule
Invalid signature      & Check 1 & Rejected \\
Tampered payload       & Check 2 & Rejected \\
Impossible sensor data & Check 5 & Rejected \\
Replay attack          & Check 4 & Rejected \\
Counter rollback       & Check 4 & Rejected \\
\midrule
\multicolumn{2}{l}{\textbf{Overall rejection rate}} & 100\% \\
\bottomrule
\end{tabular}
\end{table}

\subsection{Federated Learning Convergence}

Starting from $\mathbf{w}_1^0 = \mathbf{1}$ and
$\mathbf{w}_2^0 = \mathbf{0}$, the L2-norm difference
$\|\mathbf{w}_1^T - \mathbf{w}_2^T\|_2$ fell below 0.01 within
seven communication rounds at $\varepsilon = 50$
($\sigma = 0.097$). At $\varepsilon = 10$ ($\sigma = 0.485$),
convergence stalled at a floor of $\approx 1.37$ due to noise
domination. At $\varepsilon = 1$ ($\sigma = 4.85$), the weights
diverged. For the operational use case, $\varepsilon = 50$
provides adequate utility while maintaining a non-trivial privacy
bound on the training data.

\subsection{Scalability (20-Node Simulation)}

At 20 nodes in a $200\,\text{m} \times 200\,\text{m} $ grid,
the gossip protocol (fanout\,=\,3) delivered blocks to all nodes
within $O(\log N)$ hops; measured mean delivery latency was
[to be filled]. Packet delivery ratio at 100\,m node separation
was [to be filled] consistent with the $n=2.75$ path-loss model.
The system maintained over 80\% delivery rate with 10\% of nodes
simultaneously failed.

%% ============================================================
\section{Discussion}

Three aspects of the implementation deserve emphasis.

First, the airtime constraint is the dominant system-level
constraint, not memory or compute. At SF10/BW125/CR4-8, a
124-byte block takes 1.8\,s. This is 36\% of the 5-second
inter-block interval; when two nodes transmit simultaneously
without TDMA coordination, collision probability is high. The
TDMA scheduling in the FL layer, staggered by 1,000\,ms per
node, resolves this but limits the number of FL-capable nodes
on a shared channel to $\lfloor T_{interval} / T_{air}
\rfloor \approx 2$ without extending the round interval.
Switching to SF7 would recover the capacity at the cost of a
14\,dB link budget reduction, which may be acceptable in
short-range scenarios.

Second, the PA\_BOOST configuration issue is not unique to the
Ra-02. Any SX1276 module that routes the antenna through
PA\_BOOST rather than RFO will exhibit the same behavior with the
default library settings. Given how common this module is in
maker and academic LoRa projects, we expect this has silently
caused communication failures in published work that uses the
same hardware.

Third, on-device SHA-256 via mbedTLS has a negligible runtime
cost---1.2\,ms out of a 5,000\,ms cycle---but the include path
for mbedTLS under the ESP-IDF Arduino layer is non-obvious.
The correct header is \texttt{\#include "mbedtls/md.h"} located
in the framework's include path; adding it as a \texttt{lib\_dep}
in PlatformIO causes a duplicate symbol error.

%% ============================================================
\section{Conclusion}
\label{sec:conc}

We described a three-layer wearable monitoring system for
first-responder teams that requires no external network
infrastructure. The co-design of LoRa mesh, SHA-256 PoA
blockchain, and FedAvg with on-device differential privacy on
a single ESP32-C3 node demonstrates that all three functions can
coexist on commodity hardware costing under USD~12, within the
LoRa airtime and SRAM constraints of the platform.

The principal limitation of the current prototype is that TDMA
FL scheduling saturates with two simultaneous FL-capable nodes
on a single channel at the chosen spreading factor. Extending
to larger teams will require either a reduced spreading factor,
multiple frequency channels, or a hierarchical aggregation
topology in which cluster heads handle intra-group FL and only
aggregated updates traverse the inter-cluster LoRa links.

Hardware results are pending; we will update this paper with RSSI
measurements, actual hash-chain verification latency, and FL
convergence curves from the physical nodes.

\balance

\bibliographystyle{IEEEtran}
\bibliography{refs}

\end{document}
